% संपूर्ण मराठी ग्रंथातील प्रकरण 71: सामान्यरूपीकरण
% हा संपादनयोग्य स्रोतखंड आहे. संकलनासाठी संपूर्ण स्रोतसंग्रहातील
% अक्षररूपे, पूर्वभाग, चिन्हव्याख्या आणि आकृती-स्रोत आवश्यक आहेत.
% मूळ: Open Logic Project; मजकूर आणि रूपांतर CC BY 4.0.
% यंत्रानुवाद, दुरुस्ती आणि तपासणी: OpenAI Codex —
% GPT-5.6 Sol आणि GPT-6 Sol, दोन्ही Ultra effort.
% दुरुस्ती-पुनरावलोकन: GPT-6.1 Sol, Ultra effort.
\chapter{सामान्यरूपीकरण}\label{pt:nor:chap}










\section{परिचय}\label{pt:nor:int:sec}

~$A \lif B$ ही शर्त तुम्ही
सिद्ध केली असेल, तर
\Elim{\lif} वापरून~$B$ ची
सिद्धता करताना तिचा उपयोग
होऊ शकतो. अर्थात त्यासाठी~$A$
ची सिद्धता~$\delta_1$ ही हवीच.
~$A \lif B$ ची तुमची
सिद्धता~$\delta_3$ बहुधा
\Intro{\lif} ने संपते;
तो नियम उघड्या गृहीतक~$A$
पासून~$B$ ची
सिद्धता~$\delta_2$
घेऊन~$A \lif B$
ची सिद्धता बनवतो.
तसे असल्यास तुमची अंतिम
सिद्धता पुढील रूपाची असते:
\[
\AxiomC{$\Discharge{A}{x}$}
\RightLabel{$\delta_2$}
\DeduceC{$B$}
\DischargeRule{\Intro{\lif}}{x}
\UnaryInfC{$A \lif B$}
\AxiomC{}
\RightLabel{$\delta_1$}
\DeduceC{$A$}
\RightLabel{\Elim{\lif}}
\BinaryInfC{$B$}
\DisplayProof
\]
आधीच मिळालेली
$A \lif B$ ची सिद्धता
पुन्हा वापरली म्हणून
तुमची पद्धत कार्यक्षम असली,
तरी सिद्धता स्वतः
कार्यक्षम नाही.
$A \lif B$ आणून लगेच
निरसित करण्याचा हा
वळसा वाटतो. ~$B$
सिद्ध करण्याचा अधिक
थेट मार्ग असायला हवा.

तसा मार्ग नेहमीच असतो.
अशा “वळशांना”—परिचय
नियमाच्या लगेच नंतर
निरसन नियम लावल्यास
तयार होणाऱ्या रचनांना—
नैसर्गिक निगमनाच्या
सिद्धता मधून नेहमी
काढून टाकता येते.
याला \emph{सामान्यरूपीकरण}
म्हणतात; परिणामी मिळणारी
सिद्धता ही \emph{सामान्यरूप}
असते (किंवा ती सिद्धता
\emph{सामान्य रूपात} आहे
असे म्हणतात).

या निष्कर्षाची सिद्धता
क्रमवर्ती कलनातील
कट-निर्मूलन प्रमेयाप्रमाणेच
काहीशी गुंतागुंतीची आहे;
तिच्यात अनेक प्रसंग
पाहावे लागतात.
अंतःप्रज्ञावादी~\Log{N1i}
पेक्षा अभिजात~\Log{N1c}
साठी ती सिद्ध करणे
अधिक कठीण असते.
क्रमवर्ती कलनात~\Cut{}
नियम वापरल्यामुळे विना-कट
सिद्ध करता येत नाही असे
अधिक काही सिद्ध करता येत
नाही, हे कट-निर्मूलन दाखवते.
त्याचप्रमाणे वळसे टाळून
मिळू शकत नाहीत असे
अधिक निष्कर्ष वळसे वापरूनही
मिळत नाहीत, हे
सामान्यरूपीकरण दाखवते.
इतर साम्येही आहेत:
एखाद्या निष्कर्षाची
वळसांसह सर्वात लहान
सिद्धता जितकी असेल,
त्याच निष्कर्षाची
सामान्यरूप सिद्धता
त्याहून खूप मोठी असू शकते.
तसेच क्रमवर्ती कलनात
विना-कट सर्वात लहान
सिद्धता पेक्षा
कट असलेली सिद्धता
खूप लहान असू शकतात.
उप-सूत्र गुणधर्मही
यातून मिळतो: नैसर्गिक
निगमनाच्या सिद्धता मध्ये
निष्कर्ष आणि उघडी गृहीतके
यांच्या उप-सूत्रे, संख्यापकांमुळे लागणारी
त्यांची पद-आदेशनाची उदाहरणे आणि आवश्यक असत्य
वापरून निष्कर्ष नेहमी
सिद्ध करता येतो.
क्रमवर्ती कलनात हा गुणधर्म
कट-निर्मूलनातून मिळतो,
तसाच येथे तो
सामान्यरूपीकरणातून मिळतो.












\section{खंड आणि कट}\label{pt:nor:seg:sec}

परिचय नियमाच्या
नंतर निरसन नियम
येणे हा सिद्धता मधला
वळसा असतो; सिद्धता
सामान्यरूप करणे म्हणजे
असे वळसे काढणे.
पण~\Elim{\lor} आणि
\Elim{\lexists} मुळे
“काढायचा वळसा” याची
व्याख्याही थोडी
गुंतागुंतीची होते.
अडचण अशी: वळशात
निरसन नियम
परिचय नियमाच्या
\emph{लगेच} नंतर येईलच
असे नाही. त्यांच्यामध्ये
\Elim{\lor} किंवा
\Elim{\lexists} असू
शकतो; उदाहरणार्थ:
\[
\AxiomC{}
\DeduceC{$\lexists[x][C(x)]$}
\AxiomC{$\Discharge{A}{x}$}
\DeduceC{$B$}
\DischargeRule{\Intro{\lif}}{x}
\UnaryInfC{$A \lif B$}
\DischargeRule{\Elim{\lexists}}{y}
\BinaryInfC{$A \lif B$}
\AxiomC{}
\DeduceC{$A$}
\RightLabel{\Elim{\lif}}
\BinaryInfC{$B$}
\DisplayProof
\]
इतकेच नव्हे,~\Intro{\lif}
आणि~\Elim{\lif} यांच्या
मध्ये कितीही
\Elim{\lor} किंवा
\Elim{\lexists} अनुमाने
असू शकतात; उदाहरणार्थ:
\[
\AxiomC{}
\DeduceC{$\lexists[x][C(x)]$}
\AxiomC{}
\DeduceC{$D \lor E$}
\AxiomC{$\Discharge{A}{x}$}
\DeduceC{$B$}
\DischargeRule{\Intro{\lif}}{x}
\UnaryInfC{$A \lif B$}
\AxiomC{$\Discharge{A}{x}$}
\DeduceC{$B$}
\DischargeRule{\Intro{\lif}}{x}
\UnaryInfC{$A \lif B$}
\RightLabel{\Elim{\lor}}
\TrinaryInfC{$A \lif B$}
\DischargeRule{\Elim{\lexists}}{y}
\BinaryInfC{$A \lif B$}
\AxiomC{}
\DeduceC{$A$}
\RightLabel{\Elim{\lif}}
\BinaryInfC{$B$}
\DisplayProof
\]
या उदाहरणात तोच
\Elim{\lif} एकाहून
अधिक वळशांत भाग
घेतो: त्याच्या आधी
एकाहून अधिक~\Intro{\lif}
आहेत. अशा सर्व
शक्यता सामावण्यासाठी
~\Log{N1i}-सिद्धतेतील
सूत्र आस्थितींच्या
विशिष्ट अनुक्रमांनी
वळसे परिभाषित करतो;
त्यांना \emph{खंड}
म्हणतात. येथे त्या
~$A \lif B$
च्या आस्थितींच्या
अनुक्रमिका आहेत:
$\Elim{\lor}$ किंवा
$\Elim{\lexists}$
च्या गौण आधारविधानानंतर
त्याच अनुमानाचा
निष्कर्ष येतो.
या उदाहरणात असे
दोन खंड आहेत;
प्रत्येकात~$A \lif B$
च्या तीन आस्थिती
आहेत. आता औपचारिक
व्याख्या पाहू.

\begin{defn}
~\Log{N1i} मधील
सिद्धता~$\delta$
मध्ये एकाच सूत्र~$A$
च्या $A_1$, \dots,
$A_n$ या~$\delta$
मधील आस्थितींचा
अनुक्रम पुढील अटी
पूर्ण करत असेल,
तर तो \emph{खंड} आहे:
\begin{enumerate}
  \item $A_1$ हा
  \Elim{\lor} किंवा
  \Elim{\lexists}
  चा निष्कर्ष नसतो;
  \item $A_n$ हे
  \Elim{\lor} किंवा
  \Elim{\lexists}
  चे गौण आधारविधान नसते;
  \item $i = 1$, \dots,~$n-1$
  साठी~$A_i$ हे
  \Elim{\lor} किंवा
  \Elim{\lexists}
  चे गौण आधारविधान
  असते, आणि~$A_{i+1}$
  हा त्याच अनुमानाचा
  निष्कर्ष असतो.
\end{enumerate}
त्या खंडाची
\emph{लांबी}~$n$ आहे.
\end{defn}

प्रत्येक खंड काढून
टाकण्याजोगा वळसा
असतोच असे नाही.
अशा वळशांना
\emph{कट} म्हणतात.

\begin{defn}
एखाद्या खंडातील
$A_n$ हे
निरसन अनुमानाचे
प्रमुख आधारविधान असेल,
आणि पुढीलपैकी एक
अट पूर्ण होत असेल,
तर तो \emph{कट} आहे:
$n=1$ असून
$A_1 = A_n$ हे
परिचय अनुमानाचा
किंवा~\FalseInt चा
निष्कर्ष असेल;
किंवा~$n>1$ असेल.
\end{defn}

कट-खंडाची सुरुवात
परिचय नियमाच्या
निष्कर्षाने व्हावीच
लागते असे नाही.
तो निरसन
नियमाच्या प्रमुख
आधारविधानात संपणे
हीच आवश्यक अट आहे.
मात्र लांबी~$1$
असलेला खंड कट
म्हणून मोजायचा असेल,
तर तो परिचय
नियमाचा किंवा~\FalseInt
चा निष्कर्ष असणेही
आवश्यक आहे.

\begin{defn}
कटाची \emph{कोटी}
ही~$\depth{A}$
असते. सिद्धता~$\delta$
ची \emph{कट-कोटी}
$\cutrank{\delta}$
म्हणजे~$\delta$
मधील कटांच्या
कोटींपैकी महत्तम कोटी.
कटाची कोटी
$\cutrank{\delta}$
इतकी असेल, म्हणजे
ती महत्तम असेल,
तर तो~$\delta$
मधील \emph{महत्तम} कट.
~$\delta$ ची
\emph{कट-लांबी}
म्हणजे तिच्यातील
सर्व महत्तम कटांच्या
लांबींची बेरीज.
\end{defn}













\section{क्रमपरिवर्तन रूपांतरणे}\label{pt:nor:per:sec}

~\Log{N1i} मधील कट
हे निरसन नियमाच्या
प्रमुख आधारविधानात संपणारे
खंड असतात. सामान्यरूपीकरण
सिद्धतेतील एक प्रसंग म्हणजे
महत्तम कटांची लांबी आणि
त्याबरोबर सिद्धता ची
कट-लांबी कमी करणे.
$>1$ लांबीचा कट
अनेक प्रकारे संपू शकतो:
कटमधील~$A$ ची शेवटची
सूत्र आस्थिती
(\Elim{\lor} किंवा
\Elim{\lexists}) चा
निष्कर्ष आहे का, आणि
ती कोणत्या निरसन
अनुमानाचे प्रमुख आधारविधान
आहे, यांवर ते अवलंबून असते.

उदाहरणार्थ, संबंधित
अनुमाने~\Elim{\lor} आणि
\Elim{\land} असतील,
$A \ident B \land C$
असेल, आणि~$\delta_1$
ही उप-सिद्धता पुढील
रूपात संपत असेल:
\[
\AxiomC{}
\RightLabel{$\delta_2$}
\DeduceC{$D \lor E$}
\AxiomC{$\Discharge{D}{x}$}
\RightLabel{$\delta_3$}
\DeduceC{$B \land C$}
\AxiomC{$\Discharge{E}{y}$}
\RightLabel{$\delta_4$}
\DeduceC{$B \land C$}
\DischargeRule{\Elim{\lor}}{x\,y}
\TrinaryInfC{$B \land C$}
\RightLabel{\Elim{\land}[1]}
\UnaryInfC{$B$}
\DisplayProof
\]
तर या कटमधील शेवटची
सूत्र आस्थिती~$A_n$
ही~\Elim{\lor} चा
निष्कर्ष, म्हणजे खालचे
$B \land C$, आहे.
त्याआधीची सूत्र
आस्थिती~$A_{n-1}$
ही~\Elim{\lor}
च्या गौण आधारविधानांपैकी
एक आहे.

या कटाची लांबी
कमी करण्यासाठी
\Elim{\lor} आणि
\Elim{\land} अनुमानांचा
क्रम बदलू (“क्रमपरिवर्तन”
करू) शकतो. म्हणजे~$\delta$
मधील उपसिद्धता~$\delta_1$
पुढील रचनेने बदलतो:
\[
\AxiomC{}
\RightLabel{$\delta_2$}
\DeduceC{$D \lor E$}
\AxiomC{$\Discharge{D}{x}$}
\RightLabel{$\delta_3$}
\DeduceC{$B \land C$}
\RightLabel{\Elim{\land}[1]}
\UnaryInfC{$B$}
\AxiomC{$\Discharge{E}{y}$}
\RightLabel{$\delta_4$}
\DeduceC{$B \land C$}
\RightLabel{\Elim{\land}[1]}
\UnaryInfC{$B$}
\DischargeRule{\Elim{\lor}}{x\,y}
\TrinaryInfC{$B$}
\DisplayProof
\]
आधी~\Elim{\lor} च्या
खाली~\Elim{\land}
च्या आधारविधानात संपणारे
कट आता~\Elim{\lor}
च्या गौण आधारविधानांच्या
वर असलेल्या~\Elim{\land}
अनुमानांपैकी एका
आधारविधानात संपतात.
म्हणजे प्रत्येकाची लांबी
~$1$ ने कमी झाली.

याखेरीज~$\delta_2$,
$\delta_3$ किंवा~$\delta_4$
यांपैकी एखाद्यात पूर्णपणे
असलेला कट, किंवा~$\delta$
मधील~$\delta_1$
बाहेर पूर्णपणे असलेला कट,
बदलत नाही. त्यात तीच
सूत्रे आणि तीच
अनुमाने असतात; म्हणून
विशेषतः त्याची कोटी व
लांबी~$\delta$ मधील
अनुरूप कटाइतकीच असते.
त्यामुळे नव्या सिद्धता
ची कट-कोटी वाढत नाही,
आणि नव्या सिद्धता
ची कट-लांबी
जुन्या सिद्धता च्या
कट-लांबीपेक्षा कमी होते.

\begin{prob}
क्रमपरिवर्तन रूपांतरणानंतर
किमान एका कटाची लांबी~$1$
ने कमी होते.
\Elim{\land} ला
एखाद्या~\Elim{\lor}
च्या पलीकडे हलवले,
तर प्रत्यक्षात दोन
कट-खंडांची लांबी कमी
होते; त्यामुळे या
प्रसंगी सिद्धता ची
कट-लांबी किमान~$2$
ने कमी होते.
अशा एकाच क्रमपरिवर्तनाने
सिद्धता ची कट-लांबी
$2$ पेक्षा अधिक कमी
होऊ शकते का?
कटाची \emph{कोटी}
कमी होऊ शकते का?
\end{prob}

निरसन नियम
\Elim{\lor} किंवा
\Elim{\lexists} नियमाच्या
वर हलवण्याच्या प्रक्रियेला
\emph{क्रमपरिवर्तन रूपांतरण}
म्हणतात. काही नियम-जोड्यांचे
नमुने~\ref{pt:nor:per:tab:permutation-conversions}
मध्ये दिले आहेत.
(उरलेले सरावासाठी सोडले आहेत.)
हे नमुने वापरण्यापूर्वी मुक्ती-नियमांच्या बद्ध खुणा आणि
त्यांना बांधलेली गृहीतके एकत्र नवी नावे द्यायची,
जेणेकरून हलवलेल्या आधारसिद्धतेची उघडी गृहीतके
त्या नियमांनी अनवधानाने मुक्त होणार नाहीत.
उपसिद्धतेच्या अनेक प्रती लागल्यास प्रत्येक प्रतीतील
आयगेन स्थिरांक आणि त्यांच्याशी संबंधित बद्ध खुणा
स्वतंत्र नवी नावे देऊन पुन्हा नियमित सिद्धता मिळवायची;
उघडी गृहीतके आणि अंतिम सूत्रे जपायची.
या नावबदलांमुळे कटांची कोटी किंवा लांबी बदलत नाही.





























































































































{\small
\begin{longtable}{@{}l@{}}
\AxiomC{$D \lor E$}
\AxiomC{$\Discharge{D}{x}$}
\shortDeduce
\DeduceC{$B \land C$}
\AxiomC{$\Discharge{E}{y}$}
\shortDeduce
\DeduceC{$B \land C$}
\DischargeRule{\Elim{\lor}}{x\,y}
\TrinaryInfC{$B \land C$}
\RightLabel{\Elim{\land}[1]}
\UnaryInfC{$B$}
\DisplayProof
$\longrightarrow$\\[6ex]
\quad\AxiomC{$D \lor E$}
\AxiomC{$\Discharge{D}{x}$}
\shortDeduce
\DeduceC{$B \land C$}
\RightLabel{\Elim{\land}[1]}
\UnaryInfC{$B$}
\AxiomC{$\Discharge{E}{y}$}
\shortDeduce
\DeduceC{$B \land C$}
\RightLabel{\Elim{\land}[1]}
\UnaryInfC{$B$}
\DischargeRule{\Elim{\lor}}{x\,y}
\TrinaryInfC{$B$}
\DisplayProof
\\[10ex]
\AxiomC{$D \lor E$}
\AxiomC{$\Discharge{D}{x}$}
\shortDeduce
\DeduceC{$B \lif C$}
\AxiomC{$\Discharge{E}{y}$}
\shortDeduce
\DeduceC{$B \lif C$}
\DischargeRule{\Elim{\lor}}{x\,y}
\TrinaryInfC{$B \lif C$}
\AxiomC{$B$}
\RightLabel{\Elim{\lif}}
\BinaryInfC{$C$}
\DisplayProof
$\longrightarrow$\\[6ex]
\AxiomC{$D \lor E$}
\AxiomC{$\Discharge{D}{x}$}
\shortDeduce
\DeduceC{$B \lif C$}
\AxiomC{$B$}
\RightLabel{\Elim{\lif}}
\BinaryInfC{$C$}
\AxiomC{$\Discharge{E}{y}$}
\shortDeduce
\DeduceC{$B \lif C$}
\AxiomC{$B$}
\RightLabel{\Elim{\lif}}
\BinaryInfC{$C$}
\DischargeRule{\Elim{\lor}}{x\,y}
\TrinaryInfC{$C$}
\DisplayProof
\\[10ex]
\AxiomC{$D \lor E$}
\AxiomC{$\Discharge{D}{x}$}
\shortDeduce
\DeduceC{$B \lif C$}
\AxiomC{$\Discharge{E}{y}$}
\shortDeduce
\DeduceC{$B \lif C$}
\DischargeRule{\Elim{\lor}}{x\,y}
\TrinaryInfC{$B \lif C$}
\AxiomC{$B$}
\RightLabel{\Elim{\lif}}
\BinaryInfC{$C$}
\DisplayProof
$\longrightarrow$\\[6ex]
\AxiomC{$D \lor E$}
\AxiomC{$\Discharge{D}{x}$}
\shortDeduce
\DeduceC{$B \lif C$}
\AxiomC{$B$}
\RightLabel{\Elim{\lif}}
\BinaryInfC{$C$}
\AxiomC{$\Discharge{E}{y}$}
\shortDeduce
\DeduceC{$B \lif C$}
\AxiomC{$B$}
\RightLabel{\Elim{\lif}}
\BinaryInfC{$C$}
\DischargeRule{\Elim{\lor}}{x\,y}
\TrinaryInfC{$C$}
\DisplayProof
\\[10ex]
\AxiomC{$D \lor E$}
\AxiomC{$\Discharge{D}{x}$}
\shortDeduce
\DeduceC{$\lforall[x][B(x)]$}
\AxiomC{$\Discharge{E}{y}$}
\shortDeduce
\DeduceC{$\lforall[x][B(x)]$}
\DischargeRule{\Elim{\lor}}{x\,y}
\TrinaryInfC{$\lforall[x][B(x)]$}
\RightLabel{\Elim{\lforall}}
\UnaryInfC{$B(t)$}
\DisplayProof
\quad$\longrightarrow$\\[8ex]
\quad\AxiomC{$D \lor E$}
\AxiomC{$\Discharge{D}{x}$}
\shortDeduce
\DeduceC{$\lforall[x][B(x)]$}
\RightLabel{\Elim{\lforall}}
\UnaryInfC{$B(t)$}
\AxiomC{$\Discharge{E}{y}$}
\shortDeduce
\DeduceC{$\lforall[x][B(x)]$}
\RightLabel{\Elim{\lforall}}
\UnaryInfC{$B(t)$}
\DischargeRule{\Elim{\lor}}{x\,y}
\TrinaryInfC{$B(t)$}
\DisplayProof
\\
\caption{\Log{N1i} साठी काही क्रमपरिवर्तन रूपांतरणे.}
\label{pt:nor:per:tab:permutation-conversions}
\end{longtable}
}

\begin{prob}
\Elim{\lor} नंतर
\Elim{\lor} किंवा
\Elim{\lnot} येत असेल,
तर क्रमपरिवर्तन
रूपांतरणे द्या.
\end{prob}

\begin{prob}
\Elim{\lexists} नंतर
\Elim{\lexists} येत असेल,
तर क्रमपरिवर्तन
रूपांतरणे द्या.
दुसऱ्या प्रसंगी
आयगेन चराची कोणतीही
अट का मोडत नाही,
हे स्पष्ट करा.
\end{prob}

क्रमपरिवर्तन रूपांतरण
लावताना आयगेन चराची
अट मोडत नाही ना,
हे तपासले पाहिजे.
पुढील उदाहरण घ्या:
\[
\AxiomC{}
\RightLabel{$\delta_2$}
\DeduceC{$D \lor E$}
\AxiomC{$\Discharge{D}{x}$}
\RightLabel{$\delta_3$}
\DeduceC{$\lexists[x][B(x)]$}
\AxiomC{$\Discharge{E}{y}$}
\RightLabel{$\delta_4$}
\DeduceC{$\lexists[x][B(x)]$}
\DischargeRule{\Elim{\lor}}{x\,y}
\TrinaryInfC{$\lexists[x][B(x)]$}
\AxiomC{$\Discharge{B(c)}{z}$}
\RightLabel{$\delta_5$}
\DeduceC{$C$}
\DischargeRule{\Elim{\lexists}}{z}
\BinaryInfC{$C$}
\DisplayProof
\]
याच्या जागी पुढील
रचना बसवली, तर
\[
\AxiomC{}
\RightLabel{$\delta_2$}
\DeduceC{$D \lor E$}
\AxiomC{$\Discharge{D}{x}$}
\RightLabel{$\delta_3$}
\DeduceC{$\lexists[x][B(x)]$}
\AxiomC{$\Discharge{B(c)}{z}$}
\RightLabel{$\delta_5$}
\DeduceC{$C$}
\DischargeRule{\Elim{\lexists}}{z}
\BinaryInfC{$C$}
\AxiomC{$\Discharge{E}{y}$}
\RightLabel{$\delta_4$}
\DeduceC{$\lexists[x][B(x)]$}
\AxiomC{$\Discharge{B(c)}{z}$}
\RightLabel{$\delta_5$}
\DeduceC{$C$}
\DischargeRule{\Elim{\lexists}}{z}
\BinaryInfC{$C$}
\DischargeRule{\Elim{\lor}}{x\,y}
\TrinaryInfC{$C$}
\DisplayProof
\]
आयगेन चर~$c$ हा,
उदाहरणार्थ,~$D$
मध्ये येईल का,
अशी शंका येऊ शकते.
मूळ सिद्धता मध्ये
गृहीतक~$D$ हे
मूळ~$\Elim{\lexists}$
च्या वर मुक्त होते.
म्हणून ते
\Elim{\lexists}
अनुमानाच्या प्रमुख
आधारविधानाच्या वर
उघड्या गृहीतकात
येत नाही; त्यामुळे
तेथे आयगेन चराची अट
पूर्ण होते. पण नव्या
सिद्धता मध्ये~$D$
हे~\Elim{\lexists}
अनुमानांनंतरच मुक्त
होते, म्हणून ती अट
मोडेल असे वाटते.
तथापि, आपल्या
सिद्धता नियमित
आहेत हे लक्षात घ्या:
प्रत्येक आयगेन चर
एकाच अनुमानाचा
आयगेन चर असतो आणि
तो सिद्धता मध्ये
फक्त~\Elim{\lexists}
च्या गौण आधारविधानाच्या
वरच येतो. विशेषतः
$c$ हा~$\delta_3$
किंवा~$\delta_4$
मध्ये येऊ शकत नाही.
दाखवलेल्या दोन प्रतींतील $c$ हे एकच आयगेन स्थिरांक
दोन अनुमानांत वापरल्याचे अंतिम रूप मानायचे नाही:
वरच्या नियमानुसार प्रत्येक प्रतीला स्वतंत्र नवा स्थिरांक
आणि त्याच्या बद्ध गृहीतकांना सुसंगत नव्या खुणा द्यायच्या.

या उदाहरणातून आणखी
एक गोष्ट दिसते.
नव्या सिद्धता
मध्ये~$\delta_5$
च्या दोन प्रती आहेत.
म्हणून~$\delta_5$
मध्ये महत्तम कट असेल,
तर आपण कट-लांबी
\emph{कमी केलेली नाही}!
त्यामुळे~$\delta_5$
मध्ये महत्तम कट नाही
हे माहीत असेल तेव्हाच
क्रमपरिवर्तन रूपांतरण
लावण्याची काळजी
घ्यावी लागते.
नेहमी \emph{सर्वात वरचा}
महत्तम कट निवडून
हे साधतो: म्हणजे असा
महत्तम कट की त्याच्या
उप-सिद्धता मध्ये
(येथे~$\delta_1$)
उप-सिद्धता च्या
शेवटच्या अनुमानाच्या
वर संपणारा दुसरा
महत्तम कट नसतो.
यामुळे~$\delta_5$
मध्ये, किंबहुना
$\delta_2$, $\delta_3$
किंवा~$\delta_4$
मध्येही, इतर महत्तम
कट असण्याची शक्यता
नाहीशी होते.














\section{न्यूनीकरण रूपांतरणे}\label{pt:nor:red:sec}

~\Log{N1i} मधील
सिद्धता मध्ये एखाद्या
कटाची लांबी~$1$ असेल,
तर तो एका सूत्र
आस्थितीचा खंड असतो.
तीच सूत्र आस्थिती
परिचय अनुमानाचा किंवा~\FalseInt
चा निष्कर्ष आणि
निरसन अनुमानाचे
प्रमुख आधारविधान असते.
सामान्यरूपीकरणाच्या
सिद्धतेत अशा कटांचे
“न्यूनीकरण” करतात:
अशा परिचय/निरसन
जोडीत संपणारी उप-सिद्धता
त्या दोन्ही अनुमानांशिवाय
असलेल्या नव्या सिद्धता
ने बदलतात. याने नवे
कट निर्माण होऊ शकतात.
पण त्या नव्या कटांची
कोटी न्यूनीकृत कटाच्या
कोटीपेक्षा कमी असते.
त्यामुळे नव्या सिद्धता
ची कट-लांबी कमी होते
(कारण लांबी~$1$
असलेला महत्तम कट
निघून गेला), किंवा
सिद्धता ची कट-कोटी
कमी होते (संबंधित
महत्तम कट एकटाच
महत्तम कट असल्यास).

उदाहरणार्थ, संबंधित
अनुमाने~\Intro{\lif}
आणि~\Elim{\lif} असतील,
$A \ident B \lif C$
असेल, आणि~$\delta_1$
ही उप-सिद्धता
पुढीलप्रमाणे संपत असेल:
\[
\AxiomC{$\Discharge{B}{x}$}
\RightLabel{$\delta_2$}
\DeduceC{$C$}
\DischargeRule{\Intro{\lif}}{x}
\UnaryInfC{$B \lif C$}
\AxiomC{}
\RightLabel{$\delta_3$}
\DeduceC{$B$}
\RightLabel{\Elim{\lif}}
\BinaryInfC{$C$}
\DisplayProof
\]
हा कट काढण्यासाठी
~$\delta$ मधील
उपसिद्धता~$\delta_1$
पुढील रचनेने बदलतो:
\[
\AxiomC{}
\RightLabel{$\delta_3$}
\DeduceC{$B$}
\RightLabel{$\Subst{\delta_2}{\delta_3}{B^x}$}
\DeduceC{$C$}
\DisplayProof
\]
~$\Subst{\delta_2}{\delta_3}{B^x}$
म्हणजे~$\delta_2$
मधील~$x$ खूण असलेल्या
प्रत्येक गृहीतक~$B$
च्या जागी संपूर्ण
सिद्धता~$\delta_3$
ठेवून मिळणारी सिद्धता.
कलम-संयोजनापूर्वी $\delta_2$ च्या बद्ध मुक्ती-खुणा
आणि त्यांना बांधलेली गृहीतके एकत्र नवी नावे देऊन
$\delta_3$ च्या उघड्या खुणांपासून वेगळी ठेवायची.
त्यानंतर नव्या सिद्धतेची उघडी गृहीतके म्हणजे
$\delta_2$ मधील बसवलेल्या $B^x$ आस्थितींव्यतिरिक्त
उरलेली उघडी गृहीतके, आणि प्रत्यक्ष बसवलेल्या
$\delta_3$ च्या प्रतींची उघडी गृहीतके.
$\delta_3$ मध्ये स्वतः $B^x$ उघडे असू शकते;
म्हणून नव्या सिद्धतेत खूण $x$ अजिबात उरत नाही
असे म्हणता येत नाही. तसेच मुक्ती निष्फळ असेल किंवा
एखादी शाखा काढली गेली असेल, तर मूळची काही
उघडी गृहीतके गळू शकतात. आवश्यक निष्कर्ष असा की
नवी उघडी गृहीतके मूळ $\delta_1$ च्या उघड्या
गृहीतकांच्या पलीकडे जात नाहीत.
$\delta_2$ मधील मुक्ती-नियम आता आत बसवलेल्या
उघड्या गृहीतकांना अनवधानाने मुक्त करत नाहीत.
आणि~$\delta_3$
मधील गृहीतके मुक्त
करणारी अनुमाने अनेक
प्रतींमध्ये येऊ शकतात;
त्या प्रती
~$\Subst{\delta_2}{\delta_3}{B^x}$
मधील संबंधित गृहीतके
मुक्त करतात.
सर्व न्यूनीकरण नमुन्यांत हीच सुरक्षित कलम-संयोजनाची
अट वापरायची. अनेक प्रतींतील आयगेन स्थिरांक आणि
त्यांच्याशी संबंधित बद्ध खुणा स्वतंत्र नवी नावे देऊन
नियमितता जपायची. संख्यापकांच्या नमुन्यांतील $t$ हे
घोषित बंद पदे वापरण्याच्या संकेतानुसार बंद पद आहे.

~$\delta_2$ किंवा
$\delta_3$ यांच्या
आतमध्ये पूर्णपणे असलेला
कट, किंवा~$\delta$
मध्ये~$\delta_1$
बाहेर पूर्णपणे असलेला
कट, बदलत नाही.
त्यात तीच सूत्रे
आणि तीच अनुमाने
असतात; त्यामुळे त्याची
कोटी आणि लांबी
~$\delta$ मधील
अनुरूप कटासारखीच
राहते. आपण लांबी~$1$
असलेला महत्तम कट
काढला आहे; म्हणून
कट-लांबी कमी झाली,
किंवा~$\delta$
ची कट-लांबी~$1$
असून हा एकटाच महत्तम
कट होता, तर कट-कोटी
कमी झाली.

तथापि, दोन गोष्टी
घडू शकतात:
\begin{enumerate}
  \item सर्व~$B^x$
  गृहीतकांच्या जागी
  $\delta_3$ ठेवल्याने
  ~$\delta_3$ मधील
  सर्व कटांच्या प्रती
  वाढतात. त्यांपैकी
  काही महत्तम कट
  असतील, तर नव्या
  सिद्धता ची कट-लांबी
  जुन्या सिद्धता
  पेक्षा अधिक होऊ
  शकते. हे होणार
  नाही याची खात्री
  हवी. म्हणून न्यूनीकरण
  रूपांतरणे फक्त
  \emph{सर्वात उजव्या}
  महत्तम कटांनाच
  लावतो: न्यूनीकृत
  होत असलेल्या कटाच्या
  उजवीकडे~$\delta$
  मधील कोणत्याही
  शाखेवर दुसरा महत्तम
  कट नसतो. ~$\delta_3$
  मध्ये महत्तम कट असता, तर
  तो येथे न्यूनीकृत
  होणाऱ्या कटाच्या
  उजवीकडे असता;
  म्हणून हा सर्वात
  उजवा कट नसता.
  सर्वात उजवे कट
  निवडल्याने सिद्धता
  ची कट-लांबी किंवा
  अगदी कट-कोटीही
  कमी होते.
  \item ~$\delta_2$
  मधील एखादे गृहीतक~$B^x$
  निरसन अनुमानाचे
  प्रमुख आधारविधान
  असेल, आणि~$\delta_3$
  चे अंतिम अनुमान
  \Elim{\lor}, \Elim{\lexists}
  किंवा परिचय
  अनुमान असेल, तर
  त्या गृहीतकावर
  ~$\delta_3$ चे
  कलम-संयोजन केल्यामुळे
  नवा कट तयार होतो.
  ~$\delta_2$ मध्ये
  जे केवळ गृहीतक
  होते, ते आता लांबी~$1$
  असलेला कट
  (परिचय
  अनुमानाचा निष्कर्ष आणि
  निरसन अनुमानाचे
  प्रमुख आधारविधान)
  किंवा कट-खंडाची
  शेवटची सूत्र
  आस्थिती
  (निरसन
  अनुमानाचे प्रमुख
  आधारविधान आणि
  \Elim{\lor}
  किंवा~\Elim{\lexists}
  चा निष्कर्ष)
  होते. $B^x$ हे
  \Elim{\lor}
  किंवा~\Elim{\lexists}
  चे गौण आधारविधान
  असेल, तर ते आधीच
  एखाद्या खंडाची,
  कदाचित कट-खंडाचीही,
  पहिली सूत्र
  आस्थिती होती.
  नव्या सिद्धता
  मध्ये तो खंड
  आता अधिक लांब
  होऊ शकतो. पण
  या नव्या कटाचे
  किंवा विद्यमान
  कट-खंडाचे सूत्र
  हे~$B$ आहे, आणि
  $\depth{B} < \depth{B \lif C}$.
  म्हणून नवे कट
  आले किंवा काही
  कटांची लांबी
  वाढली, तरी त्यांपैकी
  कोणताही महत्तम कट
  नसतो. अशा रीतीने
  कट-कोटी कमी होते,
  किंवा समान कोटीचे
  कट उरले असतील
  तर कट-लांबी कमी
  होते.
\end{enumerate}

लांबी~$1$ असलेला
कट कमी करण्याच्या
प्रक्रियेला
\emph{न्यूनीकरण रूपांतरण}
(किंवा
\emph{वळसा रूपांतरण})
म्हणतात. काही
नियम-जोड्यांचे नमुने
~\ref{pt:nor:red:tab:reduction-conversions}
मध्ये दिले आहेत.
(उरलेले सरावासाठी
सोडले आहेत.)

\begin{table}
\begin{multline*}
\bottomAlignProof
\AxiomC{$\Discharge{B}{x}$}
\RightLabel{$\delta_2$}
\shortDeduce
\DeduceC{$C$}
\DischargeRule{\Intro{\lif}}{x}
\UnaryInfC{$B \lif C$}
\AxiomC{}
\RightLabel{$\delta_3$}
\shortDeduce
\DeduceC{$B$}
\RightLabel{\Elim{\lif}}
\BinaryInfC{$C$}
\DisplayProof
\quad \longrightarrow\quad
\shoveright{\bottomAlignProof
\AxiomC{}
\RightLabel{$\delta_3$}
\shortDeduce
\DeduceC{$B$}
\RightLabel{$\Subst{\delta_2}{\delta_3}{B^x}$}
\shortDeduce
\DeduceC{$C$}
\DisplayProof}
\\[4ex]
\shoveleft{\bottomAlignProof
\AxiomC{}
\RightLabel{$\delta_2$}
\shortDeduce
\DeduceC{$B$}
\AxiomC{}
\RightLabel{$\delta_3$}
\shortDeduce
\DeduceC{$C$}
\RightLabel{\Intro{\land}}
\BinaryInfC{$B \land C$}
\RightLabel{\Elim{\land}[1]}
\UnaryInfC{$B$}
\DisplayProof
\quad \longrightarrow\quad
\bottomAlignProof
\AxiomC{}
\RightLabel{$\delta_2$}
\shortDeduce
\DeduceC{$B$}
\DisplayProof}\qquad
\shoveright{\bottomAlignProof
\AxiomC{}
\RightLabel{$\delta_2$}
\shortDeduce
\DeduceC{$B$}
\AxiomC{}
\RightLabel{$\delta_3$}
\shortDeduce
\DeduceC{$C$}
\RightLabel{\Intro{\land}}
\BinaryInfC{$B \land C$}
\RightLabel{\Elim{\land}[2]}
\UnaryInfC{$C$}
\DisplayProof
\quad \longrightarrow\quad
\bottomAlignProof
\AxiomC{}
\RightLabel{$\delta_3$}
\shortDeduce
\DeduceC{$C$}
\DisplayProof}
\\[4ex]
\bottomAlignProof
\AxiomC{}
\RightLabel{$\delta_2$}
\shortDeduce
\DeduceC{$B$}
\RightLabel{\Intro{\lor}[1]}
\UnaryInfC{$B \lor C$}
\AxiomC{$\Discharge{B}{x}$}
\RightLabel{$\delta_3$}
\shortDeduce
\DeduceC{$D$}
\AxiomC{$\Discharge{C}{y}$}
\RightLabel{$\delta_4$}
\shortDeduce
\DeduceC{$D$}
\DischargeRule{\Elim{\lor}}{x\,y}
\TrinaryInfC{$D$}
\DisplayProof
\quad \longrightarrow\quad
\shoveright{\bottomAlignProof
\AxiomC{}
\RightLabel{$\delta_2$}
\shortDeduce
\DeduceC{$B$}
\RightLabel{$\Subst{\delta_3}{\delta_2}{B^x}$}
\shortDeduce
\DeduceC{$D$}
\DisplayProof}
\\[2ex]
\shoveleft{\bottomAlignProof
\AxiomC{}
\RightLabel{$\delta_2$}
\shortDeduce
\DeduceC{$B(c)$}
\RightLabel{\Intro{\lforall}}
\UnaryInfC{$\lforall[x][B(x)]$}
\RightLabel{\Elim{\lforall}}
\UnaryInfC{$B(t)$}
\DisplayProof
\quad \longrightarrow\quad
\bottomAlignProof
\AxiomC{}
\RightLabel{$\Subst{\delta_2}{t}{c}$}
\shortDeduce
\DeduceC{$B(t)$}
\DisplayProof}
\\[4ex]
\shoveright{\bottomAlignProof
\AxiomC{}
\RightLabel{$\delta_2$}
\shortDeduce
\DeduceC{$B(t)$}
\RightLabel{\Intro{\lexists}}
\UnaryInfC{$\lexists[x][B(x)]$}
\AxiomC{$\Discharge{B(c)}{x}$}
\RightLabel{$\delta_3$}
\shortDeduce
\DeduceC{$D$}
\DischargeRule{\Elim{\lexists}}{x}
\BinaryInfC{$D$}
\DisplayProof
\quad \longrightarrow\quad
\bottomAlignProof
\AxiomC{}
\RightLabel{$\delta_2$}
\shortDeduce
\DeduceC{$B(t)$}
\RightLabel{$\Subst{\Subst{\delta_3}{t}{c}}{\delta_2}{B(t)^x}$}
\shortDeduce
\DeduceC{$D$}
\DisplayProof}
\end{multline*}
\caption{\Log{N1i} साठी काही न्यूनीकरण रूपांतरणे.}
\label{pt:nor:red:tab:reduction-conversions}
\end{table}

\begin{prob}
\Intro{\lnot} नंतर
\Elim{\lnot} येत
असेल, तर न्यूनीकरण
रूपांतरण द्या.
\end{prob}

आणखी एक प्रसंग
पाहावा लागतो. कटाच्या
व्याख्येत असा प्रसंग
येऊ शकतो की कटाची
लांबी~$1$ आहे,
आणि~$A$ हे
\FalseInt{} चा निष्कर्ष
तसेच निरसन
अनुमानाचे प्रमुख
आधारविधान आहे.
$A$ हा
परिचय
नियमाचा निष्कर्ष
असताना करतो तसेच
येथेही करू. उदाहरणार्थ,
पुढील बदल करा:
\[
\bottomAlignProof
\AxiomC{}
\RightLabel{$\delta_2$}
\DeduceC{$\lfalse$}
\UnaryInfC{$B \lif C$}
\AxiomC{}
\RightLabel{$\delta_3$}
\DeduceC{$B$}
\RightLabel{\Elim{\lif}}
\BinaryInfC{$C$}
\DisplayProof
\quad\text{ऐवजी}\quad
\bottomAlignProof
\AxiomC{}
\RightLabel{$\delta_2$}
\DeduceC{$\lfalse$}
\UnaryInfC{$C$}
\DisplayProof
\]
~$A \ident (B \lor
C)$ किंवा
$A \ident \lexists[x][B(x)]$
असेल, तर थोडी काळजी
हवी. उदाहरणार्थ,
पुढील रचनेच्या जागी
\[
\bottomAlignProof
\AxiomC{}
\RightLabel{$\delta_2$}
\DeduceC{$\lfalse$}
\RightLabel{\FalseInt}
\UnaryInfC{$B \lor C$}
\AxiomC{$\Discharge{B}{x}$}
\RightLabel{$\delta_3$}
\DeduceC{$D$}
\AxiomC{$\Discharge{C}{y}$}
\RightLabel{$\delta_4$}
\DeduceC{$D$}
\DischargeRule{\Elim{\lor}}{x\,y}
\TrinaryInfC{$D$}
\DisplayProof
\quad\text{ऐवजी}\quad
\bottomAlignProof
\AxiomC{}
\RightLabel{$\delta_2$}
\DeduceC{$\lfalse$}
\RightLabel{\FalseInt}
\UnaryInfC{$D$}
\DisplayProof
\]
असा बदल केला,
तर कट-सूत्र~$D$
असलेला नवा कट
निर्माण होऊ शकतो.
पण
$\depth{D} < \depth{B \lor C}$
असेलच याची हमी नाही!{}
म्हणून~\Intro{\lor}/\Elim{\lor}
न्यूनीकरणात जसे
केले, तसेच येथेही
करावे आणि पुढील
रचनेने बदलावे:
\[
\bottomAlignProof
\AxiomC{}
\RightLabel{$\delta_2$}
\DeduceC{$\lfalse$}
\RightLabel{\FalseInt}
\UnaryInfC{$B$}
\RightLabel{$\delta_3$}
\DeduceC{$D$}
\DisplayProof
\quad\text{किंवा}\quad
\bottomAlignProof
\AxiomC{}
\RightLabel{$\delta_2$}
\DeduceC{$\lfalse$}
\RightLabel{\FalseInt}
\UnaryInfC{$C$}
\RightLabel{$\delta_4$}
\DeduceC{$D$}
\DisplayProof
\]












\section{सामान्यरूपीकरण प्रमेय}\label{pt:nor:thm:sec}

एखाद्या सिद्धता मध्ये
कट-खंड नसतील, तर ती
\emph{सामान्य} रूपात असते.
म्हणजे तिच्यावर क्रमपरिवर्तन
रूपांतरण किंवा न्यूनीकरण
रूपांतरण यांपैकी कोणतेही
लावता येत नाही; आणि
हेच या स्थितीचे लक्षण आहे.
एखादी सिद्धता
\emph{सामान्यरूप करणे}
म्हणजे कट शिल्लक
राहणार नाहीत तोपर्यंत
क्रमपरिवर्तन आणि न्यूनीकरण
रूपांतरणे लावून तिला
सामान्यरूप सिद्धतेत
बदलणे. हे नेहमी शक्य
असते, हाच
\emph{सामान्यरूपीकरण प्रमेय}.

\begin{thm}\label{pt:nor:thm:thm:normalization}
कोणतीही~\Log{N1i}-सिद्धता
$\delta$ सामान्यरूप करता
येते; म्हणजे तिचे कट
नसलेल्या सिद्धता~$\delta'$
मध्ये रूपांतर करता येते.
\end{thm}

\begin{proof}
  ~$\delta$ ची कट-कोटी
  आणि कट-लांबी यांवर
  विगमन करू. विगमन
  गृहीतक असे: कोणतीही
  सिद्धता जिची
  कट-कोटी कमी असेल,
  किंवा कट-कोटी समान
  पण कट-लांबी कमी असेल,
  सामान्यरूप करता येते.

  कट-लांबी~$0$
  असेल, तर कट नसतात
  आणि~$\delta$ आधीच
  सामान्यरूप असते.
  म्हणून~$\delta$ ची
  कट-कोटी आणि कट-लांबी
  दोन्ही~$>0$ आहेत असे
  समजू. सर्वात वरचा
  आणि सर्वात उजवा
  महत्तम कट-खंड निवडा.
  त्या कट-खंडाची लांबी
  $>1$ असेल, तर
  क्रमपरिवर्तन रूपांतरण
  लावा. लांबी~$1$
  असेल, तर संबंधित
  न्यूनीकरण रूपांतरण लावा.
  याने मिळणाऱ्या नव्या
  सिद्धता ची कट-कोटी
  तीच राहून कट-लांबी
  कमी होते, किंवा
  कट-कोटीच कमी होते;
  म्हणून विगमन गृहीतक
  लागू होते.
\end{proof}

वरील सिद्धतेतील
पद्धत—नेहमी सर्वात
वरचा आणि सर्वात
उजवा महत्तम कट
हाताळणे—सामान्यरूप
सिद्धता मिळवण्यासाठी
रूपांतरणे एका ठरावीक
क्रमाने लावावी लागतात
असे सुचवते. पण आश्चर्य
म्हणजे क्रम महत्त्वाचा
नसतो. क्रमपरिवर्तन आणि
न्यूनीकरण रूपांतरणे
कोणत्याही क्रमाने
लावली तरी शेवटी
सामान्यरूप सिद्धता
मिळते.

खंड आणि कट यांच्या
व्याख्या~\Log{N2i}
साठीही सहज देता येतात;
आणि~\Log{N1i} मधील
क्रमपरिवर्तन व न्यूनीकरण
रूपांतरणे थेट~\Log{N2i}
मध्ये रूपांतरित होतात.
त्यामुळे सामान्यरूपीकरण
प्रमेय~\Log{N2i} साठीही
खरे आहे.

\begin{thm}\label{pt:nor:thm:thm:normalization-N2i}
कोणतीही~\Log{N2i}-सिद्धता
$\delta$ सामान्यरूप करता
येते; म्हणजे तिचे कट
नसलेल्या सिद्धता~$\delta'$
मध्ये रूपांतर करता येते.
\end{thm}















\section{सामान्यरूप \Log{N2i}-सिद्धता आणि \Log{G2i} यांतील रूपांतरे}\label{pt:nor:tr:sec}

\Log{N2i} मधील सिद्धता चे \Log{G2i} मधील सिद्धता मध्ये
आणि उलट दिशेने रूपांतर करता येते. प्रथम,
कट-विरहित \Log{G2i} सिद्धता चे रूपांतर केल्यावर
सामान्यरूप \Log{N2i}-सिद्धता मिळतात.

हा निष्कर्ष मांडण्यासाठी~\Log{N2i} आणि~\Log{G2i}
यांतील क्रमवर्तींचे पूर्वपक्ष जोडणारी काही परिभाषा हवी.
खूण लावलेल्या सूत्रांचा संच~$\Gamma'$ हा खूण नसलेल्या
सूत्रांच्या बहुसंच~$\Gamma$ शी \emph{अनुरूप} असतो,
जर प्रत्येक सूत्र~$A$ साठी~$\Gamma'$ मधील
$x : A$ या रूपातील घटक ची संख्या ही~$\Gamma$
मधील~$A$ च्या प्रतींच्या संख्येहून जास्त नसेल.

\begin{cor}
$\Log{G2i} \Proves \Gamma \Sequent \Delta$ असेल, तर
$\Gamma' \Sequent C$ ची सामान्यरूप
\Log{N2i}-सिद्धता~$\delta$ असते; येथे खूण लावलेल्या
सूत्रांचा संच~$\Gamma'$ हा बहुसंच~$\Gamma$ शी अनुरूप असतो.
$\Delta=\emptyset$ असल्यास~$C\ident\lfalse$,
आणि अन्यथा~$\Delta=\{C\}$ असते.
म्हणजे, प्रत्येक सूत्र~$A$ साठी~$\Gamma'$ मधील
$x : A$ या रूपातील घटक ची संख्या ही~$\Gamma$
मधील~$A$ च्या प्रतींच्या संख्येहून (म्हणजे~$A$ च्या
पुनरावृत्तींच्या संख्येहून) जास्त नसते.
\end{cor}

\begin{proof}
\ref{pt:nat:tgi:prop:G2i-to-N2i} च्या सिद्धतेतील
रचनेकडे पाहिल्यास दिसते: परिचय नियम
सिद्धता च्या शेवटी, तर निरसन नियम
फक्त सिद्धता च्या सुरुवातीला जोडतो.
म्हणून निरसन नियमाच्या वर परिचय
नियम कधी येत नाही.
\end{proof}

मात्र~\ref{pt:nat:tgi:prop:G2i-to-N2i} मधील
\Cut{} नियमाचे रूपांतर~\Log{N2i}-सिद्धता~$\delta$ मध्ये
कट निर्माण करू शकते. एखादी~\Log{G2i} सिद्धता
शेवटी~\Cut{} ने संपत असेल, आणि तिच्या तत्काळ उप-सिद्धता
$\pi_1$ व~$\pi_2$ अनुक्रमे~$\Gamma_1 \Sequent A$
आणि~$A, \Gamma_2 \Sequent \Delta_2$ यांच्या असतील,
तर~$\pi_1$ चे रूपांतर~$\delta_1$ हे~$\pi_2$ च्या
रूपांतर~$\delta_2$ मधील $x:A \Sequent A$
स्वयंसिद्धकांवर कलम-संयोजित करतो. जर~$\delta_1$
मुख्य सूत्र~$A$ असलेल्या परिचय
नियमाने संपत असेल आणि~$\delta_2$ मधील
$x:A \Sequent A$ स्वयंसिद्धक हे निरसन
नियमाचे प्रमुख आधारविधान असेल, तर परिणामी~$\delta$
मध्ये कट तयार होतो.

\Log{N2i} मधील सिद्धता चे~\Log{G2i} मधील
सिद्धता मध्येही रूपांतर करता येते.
\ref{pt:nat:tni:prop:N2-to-G2} च्या सिद्धतेत
\Elim{\land}, \Elim{\lif}, \Elim{\lnot},
\Elim{\lforall} यांचे रूपांतर परिणामी~\Log{G2i}
सिद्धता मध्ये~\Cut{} अनुमाने निर्माण करते.
पण सुरुवातीची~\Log{N2i}-सिद्धता सामान्यरूप असेल,
तर असे कट टाळता येतात.

~\Log{N2i}-सिद्धता~$\delta$ मधील
क्रमवर्तींच्या आस्थितींच्या $S_1$, \dots,~$S_n$
या यादीला \emph{शाखा} म्हणू, जर~$S_1$ स्वयंसिद्धक
असेल आणि~$S_i$ जेव्हा एखाद्या अनुमानाचे प्रमुख
आधारविधान असेल तेव्हा~$S_{i+1}$ हा त्याचा निष्कर्ष
असेल. दुसऱ्या शब्दांत, शाखा स्वयंसिद्धकापासून
~$\delta$ मध्ये खाली जात क्रमवर्तींचा मागोवा घेते;
पण निरसन नियमांच्या गौण आधारविधानांवर थांबते.
जिचा~$S_n$ अंतिम क्रमवर्ती असतो ती~$\delta$ ची
\emph{मुख्य शाखा}. प्रत्येक सिद्धता~$\delta$ ला
किमान एक मुख्य शाखा असते, कारण प्रत्येक नियमाला
किमान एक प्रमुख आधारविधान असते (फक्त~\Intro{\land}
ला दोन असतात).

\begin{prop}\label{pt:nor:tr:prop:normal-N2i-to-G2i}
~$\delta$ ही~$\Gamma \Sequent A$ ची सामान्यरूप
$\Log{N2c}$ किंवा~$\Log{N2i}$-सिद्धता असेल,
तर~$\Gamma' \Sequent \Delta$ ची कट-विरहित
$\Log{G2c}$-सिद्धता (अनुक्रमे~$\Log{G2i}$-सिद्धता)~$\pi$
असते. येथे~$\Gamma$ मधील खुणा काढून मिळणारा
सूत्रांचा बहुसंच~$\Gamma'$ आहे;
$A$ हा~$\lfalse$ नसेल तर~$\Delta = \{A\}$,
आणि असेल तर~$\Delta = \emptyset$.
\end{prop}

\begin{proof}
यावेळी~$\Gamma \Sequent A$ च्या सिद्धता~$\delta$
मधील अनुमानांच्या संख्येवर विगमन करतो.
~$\delta$ मध्ये अनुमान नसेल, तर ती
$x:A \Sequent A$ हा स्वयंसिद्धक आहे;
तेव्हा~$\pi$ हा अनुरूप~$A \Sequent A$
स्वयंसिद्धक आहे, किंवा~$A \ident \lfalse$
असेल तर~$\lfalse \Sequent \ $ स्वयंसिद्धक आहे.

~$\delta$ चा शेवट अनुमाननियमात होत असेल,
तर पुढील प्रसंग वेगळे पाहू:
\begin{enumerate}
  \item $R$ हा परिचय नियम किंवा~\FalseInt
  असेल, तर रूपांतर~\ref{pt:nat:tni:prop:N2-to-G2}
  च्या सिद्धतेप्रमाणेच होते आणि~\Cut लागत नाही.

  \item ~$\delta$ चा शेवटचा नियम~$R$ हा निरसन नियम
  असेल, तर पुढील बाबी लक्षात घ्या:
  \begin{enumerate}
    \item ~$\delta$ ची मुख्य शाखा कोणत्याही
    परिचय नियमातून जात नाही; अन्यथा त्या
    शाखेवर परिचय नियम किंवा~\FalseInt{}
    यानंतर निरसन नियम येईल आणि~$\delta$
    सामान्यरूप राहणार नाही.
    \item ~$\delta$ च्या मुख्य शाखांवर
    परिचय नियम नसल्याने एकच मुख्य शाखा
    असते. (फक्त~\Intro{\land} ला दोन प्रमुख
    आधारविधाने आहेत; म्हणून अनेक मुख्य शाखा असतील,
    तर त्या~\Intro{\land} च्या आधारविधानांतून जाव्या लागतील.)
    \item मुख्य शाखेत~\Elim{\lor} किंवा~\Elim{\lexists}
    चे प्रमुख आधारविधान असेल, तर असा एकच नियम
    असतो आणि तो शेवटचा नियम~$R$ असतो. (नसल्यास
    शाखेत~\Elim{\lor} किंवा~\Elim{\lexists}
    चा निष्कर्ष पुढील निरसन नियमाचे प्रमुख
    आधारविधान होईल; क्रमपरिवर्तन रूपांतरण शक्य असल्यामुळे
    सिद्धता सामान्यरूप राहणार नाही.)
  \item \Elim{\lor} आणि~\Elim{\lexists} यांखेरीज
  कोणताही निरसन नियम गृहीतके निरस्त करत नाही,
  म्हणजे क्रमवर्तींचा डावा संदर्भ बदलत नाही.
  \Elim{\lor} आणि~\Elim{\lexists} हे दोन्ही नियम
  फक्त गौण आधारविधानाच्या वरची
  गृहीतके निरस्त करतात; म्हणजे फक्त गौण
  आधारविधानांचे डावे संदर्भ बदलतात.
  त्यामुळे~$\delta$ च्या मुख्य शाखेवरील
  क्रमवर्ती~$S_i$ चा संदर्भ~$\Gamma_1$ असेल,
  तर~$S_i$ हे प्रमुख आधारविधान असलेल्या अनुमानाच्या
  निष्कर्ष~$S_{i+1}$ चा संदर्भ~$\Gamma_1' \supseteq \Gamma_1$
  असतो.
  \item म्हणून मुख्य शाखेतील सर्वांत वरचा
  क्रमवर्ती~$S_1$ हा~$x: C \Sequent C$
  असेल, तर~$x:C \in \Gamma$.
  \end{enumerate}
  आता~$C$ च्या रूपानुसार प्रसंग वेगळे पाहू.
  मुख्य शाखेवरील प्रत्येक~$S_i$ हे निरसन
  नियमाचे प्रमुख आधारविधान असल्याने
  $x:C \Sequent C$ बद्दलही तेच खरे आहे.

  \item $C \ident D \land E$ असेल, तर~$\delta$
  चे रूप असे आहे:
  \[
  \Axiom$x:D \land E \fCenter D \land E$
  \RightLabel{\Elim{\land}[1]}
  \UnaryInf$x:D \land E \fCenter D$
  \Deduce$x:D \land E, \Gamma_1 \fCenter A$
  \DisplayProof
  \]
  मुख्य शाखेत~$x:D \land E \Sequent D$ आहे.
  त्या शाखेत~\Intro{\lif}, \Intro{\lnot} यांचे
  प्रमुख आधारविधान किंवा~\Elim{\lor}, \Elim{\lexists}
  यांचे गौण आधारविधान येत नाही; म्हणून
  $x:D \land E$ मधील संदर्भ-सूत्रे
  मुख्य शाखेवर बदलत नाहीत. त्यामुळे अंतिम
  क्रमवर्तीच्या संदर्भात~$x:D \land E$ असलेच पाहिजे.
  शिवाय,~\Elim{\land} ने संपणाऱ्या उप-सिद्धता
  ऐवजी~$w:D \Sequent D$ ठेवून, आणि मुख्य
  शाखेवरील प्रत्येक क्रमवर्तीच्या संदर्भात
  $x:D \land E$ ऐवजी~$w:D$ ठेवून,
  सिद्धता~$\delta$ चे सिद्धता~$\delta_1$
  मध्ये रूपांतर करता येते; म्हणजे:
  \[
  \bottomAlignProof
  \Axiom$x:D \land E \fCenter D \land E$
  \RightLabel{\Elim{\land}[1]}
  \UnaryInf$x:D \land E \fCenter D$
  \Deduce$x:D \land E, \Gamma_1 \fCenter A$
  \DisplayProof
  \quad\text{ यात }\quad
  \bottomAlignProof
  \Axiom$w:D \fCenter D$
  \Deduce$w:D, \Gamma_1 \fCenter A$
  \DisplayProof
  \]
  विगमन गृहीतक~$\delta_1$ ला लागू होते,
  कारण~$\delta$ मधील अनुमानांची संख्या ही
  $\delta_1$ मधील अनुमानांच्या संख्येपेक्षा~$1$
  ने अधिक आहे.
येथे $w$ ही संपूर्ण $\delta$ मध्ये न येणारी नवी खूण
घ्यायची. संदर्भ हे प्रत्येक नियमाच्या आधारसंदर्भांपासून
पुन्हा मोजायचे; इतर शाखेतले $x:D\land E$ किंवा
इतर गृहीतक काढायचे नाही. $\Gamma_1$ मध्ये त्या
इतर शाखांतून राहणारी सर्व गृहीतके समाविष्ट आहेत.

  विगमन गृहीतकानुसार,
  $D, \Gamma_1' \Sequent \Delta$ ची
  \Log{G2i}-सिद्धता~$\pi_1$ असते.
  \LeftR{\land} लावून~सिद्धता~$\pi$ मिळते:
  \[
  \AxiomC{}
  \RightLabel{$\pi_1$}
  \Deduce$D, \Gamma_1' \fCenter \Delta$
  \RightLabel{\LeftR{\land}}
  \UnaryInf$D \land E, \Gamma_1' \fCenter \Delta$
  \DisplayProof
  \]
  $A \ident \lfalse$ असेल, तर उत्तरपक्ष रिकामाच
  ठेवतो; उजवीकडील दुर्बलीकरण लागत नाही:
  \[
  \AxiomC{}
  \RightLabel{$\pi_1$}
  \Deduce$D, \Gamma_1' \fCenter $
  \RightLabel{\LeftR{\land}}
  \UnaryInf$D \land E, \Gamma_1' \fCenter $
  \DisplayProof
  \]
  \item $C \ident D \lor E$ असेल,
  तर ते~\Elim{\lor} चे प्रमुख आधारविधान असले
  पाहिजे आणि हे अनुमान मुख्य शाखेवरील शेवटचे
  अनुमान~$R$ असले पाहिजे. म्हणजे~$\delta$ चे
  रूप असे आहे:
  \[
  \Axiom$x:D \lor E \fCenter D \lor E$
  \AxiomC{}
  \RightLabel{$\delta_1$}
  \Deduce$y:D, \Gamma_1 \fCenter A$
  \AxiomC{}
  \RightLabel{$\delta_2$}
  \Deduce$z:E, \Gamma_2 \fCenter A$
  \RightLabel{\Elim{\lor}}
  \TrinaryInf$x:D \lor E, \Gamma_1, \Gamma_2 \fCenter A$
  \DisplayProof
  \]
  विगमन गृहीतक~सिद्धता $\delta_1$ आणि~$\delta_2$
  यांना लागू होते; त्यामुळे अनुक्रमे
  $D, \Gamma_1' \Sequent \Delta$ आणि
  $E, \Gamma_2' \Sequent \Delta$ यांच्या
  \Log{G2i}-सिद्धता $\pi_1$ व~$\pi_2$ मिळतात.
  त्यांच्यावर~$\LeftR{\lor}$ लावून~$\pi$ मिळते:
  \[
  \AxiomC{}
  \RightLabel{$\pi_1$}
  \Deduce$D, \Gamma_1' \fCenter \Delta$
  \AxiomC{}
  \RightLabel{$\pi_2$}
  \Deduce$E, \Gamma_2' \fCenter \Delta$
  \RightLabel{\LeftR{\lor}}
  \BinaryInf$D \lor E, \Gamma_1', \Gamma_2' \fCenter \Delta$
  \DisplayProof
  \]
  $\Elim{\lor}$ ने~$y: D$ किंवा~$z: E$
  निरस्त केले नसेल (म्हणजे ही गृहीतके गौण
  आधारविधानांच्या संदर्भात नसतील), तर ती आवश्यक
  \LeftR{\Weakening} ने जोडायची. $A\ident\lfalse$
  असेल, तर उत्तरपक्ष रिकामाच ठेवायचा; त्यात असत्य
  जोडायचे नाही.
  उदाहरणार्थ:
  \[
  \AxiomC{}
  \RightLabel{$\pi_1$}
  \Deduce$\Gamma_1' \fCenter $
  \RightLabel{\LeftR{\Weakening}}
  \UnaryInf$D, \Gamma_1' \fCenter $
  \AxiomC{}
  \RightLabel{$\pi_2$}
  \Deduce$\Gamma_2' \fCenter $
  \RightLabel{\LeftR{\Weakening}}
  \UnaryInf$E, \Gamma_2' \fCenter $
  \RightLabel{\LeftR{\lor}}
  \BinaryInf$D \lor E, \Gamma_1', \Gamma_2' \fCenter $
  \DisplayProof
  \]

  \item $C \ident D \lif E$ असेल,
  तर~$\delta$ चे रूप असे आहे:
  \[
  \Axiom$x:D \lif E \fCenter D \lif E$
  \AxiomC{}
  \RightLabel{$\delta_1$}
  \Deduce$\Gamma_1 \fCenter D$
  \RightLabel{\Elim{\lif}}
  \BinaryInf$x:D \lif E, \Gamma_1 \fCenter E$
  \RightLabel{$\delta_2$}
  \Deduce$x:D \lif E, \Gamma_1, \Gamma_2 \fCenter A$
  \DisplayProof
  \]
  येथे~$x:D \lif E, \Gamma_1 \Sequent E$
  समाविष्ट करणाऱ्या मुख्य शाखेत~\Intro{\lif},
  \Intro{\lnot} यांचे प्रमुख आधारविधान किंवा
  \Elim{\lor}, \Elim{\lexists} यांचे गौण
  आधारविधान नसते. त्यामुळे~$x:D \lif E, \Gamma_1$
  मधील संदर्भ-सूत्रे मुख्य शाखेवर बदलत
  नाहीत. म्हणून अंतिम क्रमवर्तीच्या संदर्भात
  $x:D \lif E, \Gamma_1$ असले पाहिजे.
  शिवाय,~\Elim{\lif} ने संपणाऱ्या उप-सिद्धता
  ऐवजी~$w:E \Sequent E$ ठेवून आणि मुख्य
  शाखेवरील प्रत्येक क्रमवर्तीच्या संदर्भात
  $x:D \lif E, \Gamma_1$ ऐवजी~$w:E$
  ठेवून,~सिद्धता च्या~$\delta_2$ या भागाचे
  योग्य सिद्धता~$\delta_2'$ मध्ये रूपांतर
  करता येते; म्हणजे:
  \[
  \bottomAlignProof
  \Axiom$x:D \lif E, \Gamma_1 \fCenter E$
  \RightLabel{$\delta_2$}
  \Deduce$x:D \lif E, \Gamma_1, \Gamma_2 \fCenter A$
  \DisplayProof
  \quad\text{ यात }\quad
  \bottomAlignProof
  \Axiom$w:E \fCenter E$
  \RightLabel{$\delta_2'$}
  \Deduce$w:E, \Gamma_2 \fCenter A$
  \DisplayProof
  \]
  विगमन गृहीतक~$\delta_1$ आणि~$\delta_2'$
  यांना लागू होते; कारण~$\delta$ मधील
  अनुमानांची संख्या ही~$\delta_1$ व~$\delta_2$
  मधील अनुमानांच्या संख्यांची बेरीज अधिक~$1$
  इतकी आहे. हा एक अधिकचा टप्पा मुख्य शाखेच्या सुरुवातीच्या
  \Elim{\lif} अनुमानासाठी आहे. (तेच अनुमान
  $\delta$ मधील शेवटचे अनुमान~$R$ असेल,
  तर~$\delta_2$ मध्ये~$0$ अनुमाने असतात.)
याही प्रसंगी $w$ ही $\delta$ मधील सर्व खुणांपासून
वेगळी नवी खूण असावी. मुख्य शाखेवरील संदर्भ नियमांनुसार
पुन्हा मोजताना इतर शाखांची गृहीतके जपायची;
$\Gamma_2$ मध्ये ती सर्व राहणारी गृहीतके येतात.
मूळ $\Gamma_1$ शी सामाईक असलेली गृहीतके नुसती
वजा करून त्या शाखांतून काढायची नाहीत.

  विगमन गृहीतकानुसार $\pi_1$ ही
  $\Gamma_1'\Sequent\Delta_D$ ची आणि $\pi_2$ ही
  $E,\Gamma_2'\Sequent\Delta$ ची \Log{G2i}-सिद्धता आहे;
  येथे $D\ident\lfalse$ असेल तर $\Delta_D=\emptyset$,
  अन्यथा $\Delta_D=\{D\}$ आहे. आवश्यक असल्यास
  पहिल्या सिद्धतेला उजवीकडील दुर्बलीकरण जोडून
  $\Gamma_1'\Sequent D$ ची $\pi_1'$ मिळते.
  \LeftR{\lif} लावून~सिद्धता~$\pi$ मिळते:
  \[
  \AxiomC{}
  \RightLabel{$\pi_1'$}
  \Deduce$\Gamma_1' \fCenter D$
  \AxiomC{}
  \RightLabel{$\pi_2$}
  \Deduce$E, \Gamma_2' \fCenter \Delta$
  \RightLabel{\LeftR{\lif}}
  \BinaryInf$D \lif E, \Gamma_1', \Gamma_2' \fCenter \Delta$
  \DisplayProof
  \]
  $D \ident \lfalse$ असेल, तर पहिल्या आधाराच्या
  रिकाम्या उत्तरपक्षात~$D$ जोडण्यासाठी एक
  \RightR{\Weakening} लावतो. $A \ident \lfalse$
  असेल, तर दुसऱ्या आधाराचा आणि अंतिम क्रमवर्तीचा
  उत्तरपक्ष रिकामाच ठेवतो. दोन्ही असत्य असताना नमुना असा:
  \[
  \AxiomC{}
  \RightLabel{$\pi_1$}
  \Deduce$\Gamma_1' \fCenter $
  \RightLabel{\RightR{\Weakening}}
  \UnaryInf$\Gamma_1' \fCenter D$
  \AxiomC{}
  \RightLabel{$\pi_2$}
  \Deduce$E, \Gamma_2' \fCenter $
  \RightLabel{\LeftR{\lif}}
  \BinaryInf$D \lif E, \Gamma_1', \Gamma_2' \fCenter $
  \DisplayProof
  \]

सामाईक खुणांची गृहीतके दोन आधारसंदर्भांत आली असल्यास
त्यांच्या खुणा काढल्यावर बहुसंचात जादा प्रती येऊ शकतात.
आवश्यक $\LeftR{\Contraction}$ लावून मूळ संचाच्या
खुणा काढून मिळणारा नेमका बहुसंच ठेवायचा.
\end{enumerate}
उरलेले प्रसंग याच प्रकारे हाताळता येतात;
ते सरावासाठी ठेवले आहेत.
\end{proof}

\begin{cor}
सामान्यरूप~\Log{N1i}-सिद्धता मध्ये येणारे प्रत्येक
सूत्र हे $\lfalse$ असते किंवा अंतिम सूत्राच्या
अथवा उघड्या गृहीतकाच्या एखाद्या उप-सूत्राचे
उदाहरण असते. सामान्यरूप~\Log{N2i}-सिद्धता साठी
अंतिम क्रमवर्तीतील सर्व सूत्रे घ्यायची.
येथे संख्यापकांमुळे उपसूत्रात मुक्त झालेल्या चलांच्या
जागी पदे बसवून मिळणारी उदाहरणेही समाविष्ट आहेत.
विधानीय प्रकरणात आदेशनाची ही अट लागत नाही.
\end{cor}

\begin{proof}
  \ref{pt:nor:tr:prop:normal-N2i-to-G2i} मधील अंतःप्रज्ञावादी
  रचना सामान्यरूप~\Log{N2i}-सिद्धता चे कट-विरहित
  \Log{G2i}-सिद्धता मध्ये रूपांतर करते.
  या रचनेत प्रत्येक असत्याव्यतिरिक्त सूत्र अनुरूप
  क्रमवर्तीच्या सिद्धतेत राहते. असत्य मात्र रिकाम्या
  उत्तरपक्षाने दर्शवले जाऊ शकते.
  कट-विरहित क्रमवर्ती नियम उलट तपासल्यास विधानीय
  आधारसूत्रे ही निष्कर्षातील सूत्रांची उपसूत्रे असतात;
  संख्यापक नियमांत त्यांच्या पद-आदेशनाची उदाहरणे येतात.
  त्यामुळे उपसूत्रांचा येथे सांगितलेला गुणधर्म मिळतो.
  \Log{N1i} चे~\Log{N2i} मध्ये आधीचे रूपांतर सूत्रे
  आणि उघडी गृहीतके जपते, म्हणून त्यालाही हा निष्कर्ष लागू होतो.
\end{proof}



